\clearpage
\begin{proofbox}

    \textbf{\textcolor{CProof}{证明前先看旧知识：为什么“平移”会把一个值整个搬走？}}

    从 $y=\dfrac{k}{x}$（$k\ne0$）出发。
    原图上任何一个点 $(u,k/u)$，横坐标、纵坐标都不为 $0$。
    若整张图先向右移 $h$，再向上移 $v$，这个点就变成
    \[
       \left(u+h,\,\frac{k}{u}+v\right).
    \]
    因为 $u\ne0$，新点横坐标 $u+h\ne h$；因为 $k/u\ne0$，新点纵坐标 $k/u+v\ne v$。
    反过来，给定任意 $t\ne v$，取 $u=\dfrac{k}{t-v}$，就能得到纵坐标为 $t$ 的新点。
    因此平移后的函数\textbf{不可能取到 $v$，但能取到其余每个实数}。
    这就从初中学过的反比例函数出发，直接解释了为什么最后的值域“只少一个数”。

    \medskip
    \textbf{如何接上一般的一次分式？}
    把 $\dfrac{ax+b}{cx+d}$ 拆成
    \[
       y=\underbrace{\frac ac}_{v}+
          \frac{\overbrace{(bc-ad)/c^2}^{k}}{x-\underbrace{(-d/c)}_{h}}.
    \]
    由于 $k\ne0$，上面的平移解释可以直接使用，于是取不到的纵坐标是 $v=a/c$。
    下面的三个正式证明，再分别用凑分母、反解 $x$、图像与代数运算把结论算清楚。

    \medskip
    \textbf{\textcolor{CProof}{证明一：凑出一个分母（最容易理解，推荐先学）}}
\begin{center}\begin{tikzpicture}[font=\small,>=stealth,line join=round]
\draw[gray!45,rounded corners=3pt] (0,0) rectangle (14.2,7.2);
\node at (7.1,6.6) {分离常数后，看清两个永远取不到的坐标};
\begin{scope}[shift={(3.8,2.8)},scale=.58]
\draw[->] (-3.2,0)--(4.4,0) node[right]{$x$};\draw[->] (0,-2.1)--(0,5.4) node[above]{$y$};
\draw[dashed,magenta] (1,-2)--(1,5.1);\draw[dashed,magenta] (-3.1,2)--(4.2,2);
\draw[teal!80!black,thick,domain=-3:.22,samples=65] plot(\x,{2+3/(\x-1)});
\draw[teal!80!black,thick,domain=2:4.1,samples=65] plot(\x,{2+3/(\x-1)});
\draw (1,2) circle(2pt);
\node[below] at (1,0){$1$};\node[above left] at (0,2){$2$};

\end{scope}
\node[align=left,anchor=west] at (8.1,3.9){$x\ne1$，$y\ne2$\\[10pt] 中心 $(1,2)$\\[10pt] $(x-1)(y-2)=3$};
\node[text width=13.2cm,align=center] at (7.1,.65) {示例 $y=(2x+1)/(x-1)=2+3/(x-1)$；虚线为两条渐近线。};
\end{tikzpicture}
\end{center}


    \textbf{要证什么：} 对 $y=\dfrac{ax+b}{cx+d}$，当 $c\ne0$、$ad\ne bc$ 且 $x\ne-d/c$ 时，证明函数值能取遍除 $a/c$ 外的所有实数。

    \begin{description}[style=nextline, leftmargin=2.8em, labelsep=0.5em, itemsep=0.25em]
        \item[\textbf{第 1 步：先想小学的“除法”。}]
            “被除的数 $=$ 除数 $\times$ 商 $+$ 剩下的数”。这里想把 $ax+b$ 写成“$cx+d$ 的倍数 $+$ 剩下的数”。
            $ax$ 要由 $cx$ 凑出来，倍数必须是 $a/c$，因此
            \[
                ax+b=\frac ac(cx+d)+\frac{bc-ad}{c}.
            \]
        \item[\textbf{第 2 步：两边都除以分母。}]
            因为 $cx+d\ne0$，所以
            \[
                y=\frac ac+\frac{bc-ad}{c(cx+d)}.
            \]
        \item[\textbf{第 3 步：先找出取不到的数。}]
            $ad\ne bc$ 说明 $bc-ad\ne0$。一个非零的数除以另一个非零的数，结果不会是 $0$，因此 $y\ne a/c$。
        \item[\textbf{第 4 步：再证明其他数都能取到。}]
            任取 $t\ne a/c$，希望 $y=t$。由上式只需让
            $\dfrac{bc-ad}{c(cx+d)}=t-\dfrac ac$。
            解这个等式得到
            \[
                x=-\frac dc+\frac{bc-ad}{c^2\left(t-a/c\right)}.
            \]
            因为 $t-a/c\ne0$、$bc-ad\ne0$，所以这个 $x$ 存在，而且 $x\ne-d/c$。于是 $y=t$ 一定能够实现。
    \end{description}
    \textbf{得到：} $y$ 恰好不能等于 $a/c$，除此之外没有别的遗漏。

    \medskip
    \end{proofbox}
\clearpage
\begin{proofbox}
\textbf{\textcolor{CProof}{证明二：把 $y$ 当成已知数，直接解 $x$（考试最快）}}
\begin{center}\begin{tikzpicture}[font=\small,>=stealth,line join=round]
\draw[gray!45,rounded corners=3pt] (0,0) rectangle (14.2,9.4);
\node at (7.1,8.8) {给定输出值，看水平线能否与图像相交};
\begin{scope}[shift={(3.8,3.4)},scale=.8]
\draw[->] (-3.2,0)--(4.4,0) node[right]{$x$};\draw[->] (0,-2.1)--(0,5.4) node[above]{$y$};
\draw[dashed,magenta] (1,-2)--(1,5.1);\draw[dashed,magenta] (-3.1,2)--(4.2,2);
\draw[teal!80!black,thick,domain=-3:.22,samples=65] plot(\x,{2+3/(\x-1)});
\draw[teal!80!black,thick,domain=2:4.1,samples=65] plot(\x,{2+3/(\x-1)});
\draw (1,2) circle(2pt);
\node[below] at (1,0){$1$};\node[above left] at (0,2){$2$};
\draw[blue,thick] (-3,5)--(4.1,5);\fill[blue] (2,5) circle(2pt);\draw[dashed] (2,0)--(2,5);\node[below] at (2,0){$2$};\node[above right] at (2,5){$(2,5)$};
\end{scope}\node[align=left,anchor=west] at (8.6,5.1){$t=5$：$x=2$\\[9pt] $t=2$：没有交点\\[12pt] $t\ne2$ 时\\[4pt] $x=1+\dfrac3{t-2}$};
\node[text width=13.2cm,align=center] at (7.1,.65) {水平线 $y=2$ 不相交；每条 $y=t\ne2$ 都有一个合法交点。};
\end{tikzpicture}
\end{center}


    \begin{description}[style=nextline, leftmargin=2.8em, labelsep=0.5em, itemsep=0.25em]
        \item[\textbf{第 1 步：设想一个想要得到的值 $t$。}]
            要让函数值等于 $t$，必须有
            \[
                t=\frac{ax+b}{cx+d}\quad\Longleftrightarrow\quad
                t(cx+d)=ax+b.
            \]
        \item[\textbf{第 2 步：把含 $x$ 的项放在一起。}]
            展开、移项：
            \[
                (ct-a)x=b-dt.
            \]
        \item[\textbf{第 3 步：分情况。}]
            如果 $t=a/c$，左边的系数 $ct-a=0$，而右边
            \[
                b-d\frac ac=\frac{bc-ad}{c}\ne0.
            \]
            这变成了 $0\cdot x=\text{非零的数}$，不可能成立。
            如果 $t\ne a/c$，则
            \[
                x=\frac{b-dt}{ct-a}
            \]
            是唯一的解。并且
            \[
                cx+d=\frac{bc-ad}{ct-a}\ne0,
            \]
            所以它没有违反分母不能为零的限制。
    \end{description}
    \textbf{得到：} $t=a/c$ 不可能，其余每一个 $t$ 都有对应的合法 $x$。

    \medskip
    \end{proofbox}
\clearpage
\begin{proofbox}
\textbf{\textcolor{CProof}{证明三：化成初中学过的反比例函数（图像思路）}}
\begin{center}\begin{tikzpicture}[font=\small,>=stealth,line join=round]
\draw[gray!45,rounded corners=3pt] (0,0) rectangle (14.2,7.2);
\node at (7.1,6.6) {平移坐标后，两条曲线就是同一个反比例关系};

\begin{scope}[shift={(3.7,3.6)},scale=.67]
\draw[->] (-3.2,0)--(3.4,0) node[right]{$X$};\draw[->] (0,-3.1)--(0,3.2) node[above]{$Y$};
\draw[teal!80!black,thick,domain=1:3.1] plot(\x,{3/\x});\draw[teal!80!black,thick,domain=-3.1:-1] plot(\x,{3/\x});
\draw[dashed,gray] (-1.5,-2)--(1.5,2);\fill[blue] (1.5,2) circle(2pt);\fill[orange!90!black] (-1.5,-2) circle(2pt);\node[above right] at (1.5,2){$(X,Y)$};\node[below left] at (-1.5,-2){$(-X,-Y)$};\node[below right] at (0,0){$O'$};
\end{scope}
\node[align=left,anchor=west] at (8.2,3.9){$X=x-1,\ Y=y-2$\\[9pt] $XY=3$\\[9pt] $Y=3/X\ne0$\\[9pt] 平移后中心为 $(1,2)$};
\node[text width=13.2cm,align=center] at (7.1,.65) {一般情形：$X=x+d/c$，$Y=y-a/c$，$XY=(bc-ad)/c^2\ne0$。};
\end{tikzpicture}
\end{center}


    \begin{description}[style=nextline, leftmargin=2.8em, labelsep=0.5em, itemsep=0.25em]
        \item[\textbf{第 1 步：把式子改写成两个括号相乘。}]
            由证明一的等式，移项并整理：
            \[
                \left(x+\frac dc\right)\left(y-\frac ac\right)
                =\frac{bc-ad}{c^2}.
            \]
        \item[\textbf{第 2 步：给两个括号取个简短名字。}]
            设 $X=x+d/c$、$Y=y-a/c$，并记 $k=(bc-ad)/c^2$。由于 $ad\ne bc$，$k\ne0$，于是
            \[
                XY=k,\qquad Y=\frac{k}{X},\qquad X\ne0.
            \]
        \item[\textbf{第 3 步：用反比例函数的旧知识。}]
            $k/X$ 永远不为 $0$。反过来，如果想让它等于任意一个非零数 $q$，只需取 $X=k/q$，这是一个合法的非零数。因此 $Y$ 可以是任何非零实数。
        \item[\textbf{第 4 步：把名字换回去。}]
            $Y=y-a/c\ne0$，而且它能取遍所有非零实数。因此 $y$ 能取遍除 $a/c$ 外的所有实数。
    \end{description}

    \medskip
    \textbf{\textcolor{CProof}{补充：为什么图像中心也是这个点？}}
    上面得到 $XY=k$。如果 $(X,Y)$ 满足它，那么 $(-X,-Y)$ 也满足，因为 $(-X)(-Y)=XY=k$。这两个点恰好以 $(0,0)$ 为中心对称。
    换回原来的 $x,y$，中心就是 $\left(-d/c,\,a/c\right)$。

    \medskip

    \medskip
    \textbf{\textcolor{CProof}{再向前一步：不能只盯着“中心不在曲线上”。}}
    原反比例函数的对称中心 $(0,0)$ 不在图像上；更重要的是，
    \textbf{$y=0$ 这条横线上的任何点都不在图像上}。
    平移后，中心变成 $(-d/c,\,a/c)$，而整条没有图像点的横线变成 $y=a/c$，
    所以 $a/c$ 才是取不到的函数值。\textbf{仅凭某一个点不在图像上，不能推出它的纵坐标取不到。}

    \textbf{\textcolor{CProof}{补充：不用求导，怎样看增减？}}
    在 $x=-d/c$ 同一侧取 $x_1<x_2$。直接通分可得
    \[
        \begin{aligned}
        y(x_2)-y(x_1)
          &=\frac{ax_2+b}{cx_2+d}-\frac{ax_1+b}{cx_1+d}\\
          &=\frac{(ad-bc)(x_2-x_1)}{(cx_2+d)(cx_1+d)}.
        \end{aligned}
    \]
    因为 $x_1,x_2$ 在 $-d/c$ 同一侧，两个分母同号，乘积为正；又因 $x_2-x_1>0$，所以差的正负只看 $ad-bc$：正则递增，负则递减。
    \textbf{不能把两侧合起来当成一个区间判断增减。}
\end{proofbox}
